%! TEX root = ../m1-19-20.tex

\chapter{Convergența seriilor --- Exerciții}

Studiați convergența seriilor de forma $ \sum x_n $. În cazul
seriilor alternante, decideți și convergența absolută sau semiconvergența:
\begin{enumerate}[(1)]
    \item $ x_n = (\arctan 1)^n $; \hfill(D, radical)
    \item $ x_n = \sqrt{n} \cdot \ln \left( 1 + \dfrac{1}{n} \right) $; \hfill(D, comparație la limită)
    \item $ x_n = \dfrac{1}{n - \ln n} $; \hfill(D, comparație la limită)
    \item $ x_n = (-1)^n \dfrac{n + 1}{n^3} $; \hfill(C, Leibniz)
    \item $ x_n = \dfrac{\ln n}{2n^3 - 1} $; \hfill(C, comparație)
    \item $ x_n = \dfrac{\sqrt[n]{2}}{n^2} $; \hfill(C, comparație/integral)
    \item $ x_n = \dfrac{\sqrt{n + 1} - \sqrt{n}}{\sqrt[3]{n^2}} $; \hfill(C, comparație)
    \item $ x_n = \dfrac{1}{n \ln^2 n} $; \hfill(C, integral)
    \item $ x_n = \dfrac{\ln n}{n^2} $; \hfill(C, comparație)
    \item $ x_n = \dfrac{e^\frac{1}{n}}{n} $; \hfill(D, comparație la limită)
    \item $ x_n = \dfrac{(-1)^n \ln n}{\sqrt{n}} $; \hfill(C, Leibniz)
    \item $ x_n = \dfrac{(-1)^n n!}{(2n)!} $; \hfill(C, Leibniz)
    \item $ x_n = \dfrac{\arctan n}{n^2 + 1} $; \hfill(C, comparație/integral)
    \item $ x_n = \dfrac{\sqrt{n}}{\ln(n + 1)} $; \hfill(D, necesar)
    \item $ x_n = (-1)^{n+1} \dfrac{2n + 1}{3^n} $; \hfill(C, Leibniz)
    \item $ x_n = \dfrac{\sqrt[3]{n^3 + n^2} - n}{n^2} $; \hfill(parțial 2018--2019)
    \item $ x_n = (-1)^n \dfrac{n}{n^2 - 1} $; \hfill(parțial 2018--2019)
    \item $ x_n = (-1)^n (\sqrt{n^2 + 1} - n) $; \hfill(parțial 2018--2019)
    \item $ x_n = \dfrac{a^n}{2n^2 + 1} $; \hfill(discuție după $ a $)
    \item $ x_n = (-1)^n \dfrac{(2n)!}{n^n} $; \hfill(parțial FIA)
    \item $ x_n = \left( \dfrac{n}{n + a} \right)^{n^2}, a > 0 $; \hfill(discuție după $ a $)
    \item $ x_n = \dfrac{(-1)^n}{2^{2n} \cdot n!} $ \hfill(parțial FIA)
    \item $ x_n = \dfrac{a^n}{n^3}, a > 0 $. \hfill(discuție după $ a $)
\end{enumerate}

%%% Local Variables:
%%% mode: latex
%%% TeX-master: "../m1-19-20"
%%% End:
